\section{Technical comparison of solution routes}
\label{sec:comparison}

\begin{landscape}
\scriptsize\sloppy
\renewcommand{\arraystretch}{0.88}
\setlength{\tabcolsep}{2.5pt}
\begin{longtable}{@{}>{\raggedright\arraybackslash}p{.13\linewidth}>{\raggedright\arraybackslash}p{.18\linewidth}>{\raggedright\arraybackslash}p{.15\linewidth}>{\raggedright\arraybackslash}p{.20\linewidth}>{\raggedright\arraybackslash}p{.22\linewidth}@{}}
\caption{What each method simplifies and what it costs.}\label{tab:comparison}\\
\toprule
Method & Main simplification & Remaining solve & Embedded assessment & Nonlinear extension / didactic value\\
\midrule
\endfirsthead
\toprule
Method & Main simplification & Remaining solve & Embedded assessment & Nonlinear extension / didactic value\\
\midrule
\endhead
Direct Lagrange & One smooth active equality; parallel normals & Coupled polynomial candidates & Exact candidates, but branch and global checks required & Useful local check; teaches tangency after the picture is known\\
KKT active set & Multiple active boundaries and status logic & Enumerate active sets and roots & Robust verification language; not automatically a global solver & Natural for local affine models and added bounds\\
Generalized eigenvalue & Torque per fixed loss when voltage is inactive & One extreme eigenpair & Excellent: small, symmetric, well conditioned after whitening & Recompute locally; strong resource-efficiency interpretation\\
Whitening plus torque rotation & Loss sphere and \(M=\kappa z_\psi z_q\) & Generic voltage quadratic in 3D & Very cheap transforms; ideal preprocessing & Local metrics remain useful; strongest physical coordinate view\\
Active flux & Isolates torque-producing flux & Needs independent loss-orthogonal complement & Excellent physical signal, but not alone a complete optimizer & Established machine concept; intuitive bridge\\
Hyperbolic coordinates & Constant torque becomes fixed \(i_T\) & Allocate \(\alpha,z_0\) under voltage & Transcendentals unnecessary if replaced by ratios & Excellent explanation of imbalance; moderate solver value\\
Projective ratios & Removes radius; loss levels become circles & Tangency of two conics, generic sextic stationary equation & Good 2D scalar/curve solver; explicit branch \(\eta>0\) & Local affine voltage still conic-like; excellent visual value\\
Complex square map & Hyperbolas become lines; loss radius retained & Voltage introduces roots and branches, especially with \(z_0\) & Square-root inverse is cheap, but full constraint is not simpler & Elegant interpretation, weak primary-solver choice\\
Resultants / discriminants & Removes variables; contact is repeated root & Quartic intersections; up to sextic stationary polynomial here & Exact offline/reference route; root conditioning and selection costly & PWA region by region; shows why tangency changes root multiplicity\\
Performance space & Direction maps to convex \((\nu,\tau)\) set & One support slope and \(3\times3\) extreme eigenpairs & Preferred: bounded scalar search, no trig, explicit status & Rebuild locally; best unification of both operating problems\\
SVD / GSVD & Voltage per loss principal gains and null direction & Torque remains coupled & Excellent diagnostics and scaling; no universal final basis & Local Jacobian SVD is valuable; clarifies conditioning\\
PMSM quartics & Exact intersections of 2D quadrics & Fourth-order roots and physical filtering & Viable with careful libraries; often less robust than bracket search & Map-updated coefficients possible; mature reference literature\\
Fixed-point outer loop & Preserves an analytic inner solve & A few map updates and residual tests & Promising if contraction/damping is monitored & Direct route to saturation and cross-saturation\\
\bottomrule
\end{longtable}
\end{landscape}

No method wins because it is exotic. The most useful combination is mundane
and structured: whitening for scaling, active-flux rotation for understanding,
performance-space support for online direction selection, and KKT/original
equation residuals for verification. Projective circles are the clearest 2D
picture. The square map and stereographic projection are mathematically
beautiful but do not sufficiently simplify the full voltage constraint.

For an \gls{fpga}, a small fixed-iteration eigensolver or tabled support curve
offers regular data flow. For an \gls{mcu}, the same boundary can be found with
bisection and a compact symmetric eigensolver. A radical quartic is exact in a
symbolic sense, yet can be the less predictable numerical implementation once
coefficient scaling, nearly repeated roots, and physical branch filtering are
counted.
